\documentclass[11pt]{article}

\usepackage[margin=1.05in]{geometry}
\usepackage[T1]{fontenc}
\usepackage{lmodern}
\usepackage{microtype}
\usepackage{amsmath,amssymb,amsthm,mathtools}
\usepackage{booktabs}
\usepackage{array}
\usepackage{enumitem}
\usepackage{xcolor}
\usepackage{listings}
\usepackage[hidelinks]{hyperref}
\usepackage[nameinlink,noabbrev]{cleveref}
\setlength{\emergencystretch}{2em}

\definecolor{codeblue}{RGB}{31,78,121}
\definecolor{codegray}{RGB}{82,82,82}
\lstdefinelanguage{Lean}{
  morekeywords={abbrev,def,theorem,lemma,namespace,end,import,where,if,then,else,
    let,in,match,with,fun,forall,exists,by,exact,refine,have,show,from,calc,
    unfold,rw,simp,simpa,apply,intro,rintro,constructor},
  sensitive=true,
  morecomment=[l]{--},
  morecomment=[s]{/-}{-/},
  morestring=[b]"
}
\lstset{
  basicstyle=\ttfamily\footnotesize,
  keywordstyle=\color{codeblue}\bfseries,
  commentstyle=\color{codegray}\itshape,
  breaklines=true,
  columns=fullflexible,
  frame=single,
  rulecolor=\color{black!15},
  xleftmargin=0.4em,
  xrightmargin=0.4em,
  literate={∧}{{$\land$}}1 {¬}{{$\neg$}}1 {∃}{{$\exists$}}1 {→}{{$\to$}}1
}

\newtheorem{theorem}{Theorem}[section]
\newtheorem{proposition}[theorem]{Proposition}
\newtheorem{lemma}[theorem]{Lemma}
\newtheorem{corollary}[theorem]{Corollary}
\theoremstyle{definition}
\newtheorem{definition}[theorem]{Definition}
\newtheorem{remark}[theorem]{Remark}

\newcommand{\bits}{\{0,1\}}
\newcommand{\D}{\mathcal{D}}
\newcommand{\C}{\mathcal{C}}
\newcommand{\K}{\mathcal{K}}
\newcommand{\lcs}{\operatorname{LCS}}

\title{\bfseries The Largest Binary Four-Deletion-Correcting\\
Code of Length Twelve}
\author{Anonymous}
\date{September 28, 2026}

\begin{document}
\maketitle

\begin{abstract}
A binary code of length $n$ corrects $t$ deletions when no two codewords
share a subsequence of length $n-t$.  Let $L_2(n,t)$ be the largest size of
such a code.  For length twelve and four deletions, Kolev's 2016 table
records $6\leq L_2(12,4)\leq10$, and a 2008 table of variable-length
insertion/deletion codes by Baker, Flack, and Houghten implies
$L_2(12,4)\leq7$.  We prove
\[
  L_2(12,4)=6.
\]
A six-word code gives the lower bound.  For the upper bound we construct an
exact weighted-packing certificate: 117 nonzero weights on length-eight
descendant rows and 20 weighted cliques of the conflict graph.  Each of the
4,096 binary words of length twelve receives weight at least one, and the
weights sum to $6.999616<7$, so no code has seven words.  A standard-library
Python program checks the certificate in integer arithmetic.  A Lean~4
development proves the weighted-sum argument and checks the certificate data.
\end{abstract}

\section{Introduction}

A deletion channel drops symbols and does not tell the receiver where.
Levenshtein founded the theory of codes that survive such errors
\cite{levenshtein1966binary}.  A deletion shifts every later symbol, so the
Hamming-distance tools of classical coding theory do not apply, and the exact
extremal numbers remain unknown even for short binary words.

The finite values of $L_2(n,t)$ have a short history.  Landjev and Haralambiev
tabulated lower and upper bounds for $t=2,3,4,5$ in 2007
\cite{landjev2007multiple}; at $(n,t)=(12,4)$ their tables give
$5\leq L_2(12,4)\leq10$.  Kolev's 2016 table raises the lower end to six and
credits that value to a 2015 Sofia University master's thesis
\cite[Table~1]{kolev2016optimal}, \cite{landgev2015thesis}.  A sharper upper bound follows from a
neighbouring problem.  Baker, Flack, and Houghten computed optimal
variable-length insertion/deletion codes, in which $n$ is the length of the
longest codeword, and their table lists $D_2(12,10)=7$
\cite{baker2008variable,houghtenTable}.  Two words of length twelve are at
insertion/deletion distance $24-2\lcs(x,y)$, so a four-deletion-correcting
code of length twelve has minimum distance at least ten and counts as a
variable-length code.  Hence $L_2(12,4)\leq7$, although no fixed-length table
records this.  Recent computational work concentrates on
fewer deletions: Nakasho, Hagiwara, Anderson, and Nation proved
$L_2(11,1)=172$ with integer programming \cite{nakasho2023tight}, and Weindel
and Heckel tabulate bounds for two and three deletions up to length eighteen
\cite{weindel2026llm}.  Neither work treats four deletions.  \Cref{tab:history} summarizes the record at $(12,4)$.

\begin{table}[h]
\centering
\begin{tabular}{@{}lcc@{}}
\toprule
Source & Lower bound & Upper bound\\
\midrule
Landjev and Haralambiev (2007) \cite{landjev2007multiple} & 5 & 10\\
Baker, Flack, and Houghten (2008), via $D_2(12,10)$ \cite{baker2008variable} & & 7\\
Kolev (2016) \cite{kolev2016optimal} & 6 & 10\\
This paper & 6 & \textbf{6}\\
\bottomrule
\end{tabular}
\caption{Published bounds on $L_2(12,4)$.}
\label{tab:history}
\end{table}

\begin{theorem}[Main theorem]\label{thm:intro}
The largest binary code of length twelve that corrects four deletions has six
codewords.  Equivalently, among any seven binary words of length twelve, some
two share a common subsequence of length eight.
\end{theorem}

The six-word code in \cref{sec:lower} reproduces the known lower bound.  The
new content is the upper bound: seven codewords cannot coexist, which improves
the implied bound of seven by one.
The argument has three parts.
\begin{enumerate}[leftmargin=2.2em]
  \item A weighted-packing lemma (\cref{lem:weighted}) combines the standard
        descendant-row constraints with inequalities from cliques of the
        conflict graph.
  \item An explicit certificate (\cref{sec:certificate}) supplies
        nonnegative integer weights whose normalized total is $6.999616$.
  \item Two independent checkers, a standard-library Python program and a
        Lean~4 development, verify every coefficient, coverage inequality,
        and clique pair.
\end{enumerate}

We found the certificate with floating-point linear programming, but the
theorem depends only on the exported integers and one weighted sum.  The
proof inherits nothing from the solver.

\section{Definitions and the conflict hypergraph}

Let $[q]=\{0,1,\dots,q-1\}$.  For $x=x_1x_2\cdots x_n\in[q]^n$ and
$0\leq t<n$, define
\[
  \D_t(x)
  =\bigl\{x_{i_1}x_{i_2}\cdots x_{i_{n-t}}:
      1\leq i_1<i_2<\cdots<i_{n-t}\leq n\bigr\}.
\]
Different choices of deleted positions can produce the same word; that word
counts once in $\D_t(x)$.

\begin{definition}
A code $\C\subseteq[q]^n$ corrects $t$ deletions if
\[
  \D_t(x)\cap\D_t(y)=\varnothing
  \qquad\text{for all distinct }x,y\in\C.
\]
Let $L_q(n,t)$ denote the maximum cardinality of such a code.
\end{definition}

Some authors write $A_q(n,t)$ for the same quantity.  We follow the notation
of the finite tables in \cite{landjev2007multiple,kolev2016optimal}.

Fix $(q,n,t)$ and form a hypergraph on the vertex set $[q]^n$.  Each
$d\in[q]^{n-t}$ indexes a row
\[
  R_d=\{x\in[q]^n:d\in\D_t(x)\}.
\]
A deletion-correcting code meets each $R_d$ in at most one vertex.  This
set-packing formulation underlies the non-asymptotic upper bounds of Kulkarni
and Kiyavash \cite{kulkarni2013nonasymptotic}.

The \emph{conflict graph} $G$ has vertex set $[q]^n$, with two distinct words
adjacent when they share a deletion descendant.  Codes correcting $t$
deletions are the independent sets of $G$, so a code meets any clique
$K\subseteq V(G)$ in at most one word.

For equal-length words $x,y$, let $\lcs(x,y)$ be the length of a longest
common subsequence.

\begin{lemma}\label{lem:lcs}
For $x,y\in[q]^n$, the sets $\D_t(x)$ and $\D_t(y)$ intersect if and only if
$\lcs(x,y)\geq n-t$.
\end{lemma}

\begin{proof}
A common descendant is a common subsequence of length $n-t$.  Conversely, a
common subsequence of length at least $n-t$ contains one of length exactly
$n-t$, and that word lies in both descendant sets.
\end{proof}

\section{A code of size six}\label{sec:lower}

Consider the six binary words
\begin{equation}\label{eq:codebook}
\begin{split}
\C_0=\{&000000000000,\quad 000011111110,\quad 010101010101,\\
       &101111100000,\quad 110000000111,\quad 111111111111\}.
\end{split}
\end{equation}
Their pairwise longest-common-subsequence lengths, in the order listed, are:

\begin{center}
\begin{tabular}{c@{\qquad}rrrrrr}
\toprule
 & 1 & 2 & 3 & 4 & 5 & 6\\
\midrule
1 & 12 & 5 & 6 & 6 & 7 & 0\\
2 & 5 & 12 & 7 & 7 & 7 & 7\\
3 & 6 & 7 & 12 & 7 & 7 & 6\\
4 & 6 & 7 & 7 & 12 & 7 & 6\\
5 & 7 & 7 & 7 & 7 & 12 & 5\\
6 & 0 & 7 & 6 & 6 & 5 & 12\\
\bottomrule
\end{tabular}
\end{center}

Each off-diagonal entry is at most $7<12-4$, so by \cref{lem:lcs} the
four-deletion descendant sets are pairwise disjoint.

\begin{proposition}\label{prop:lower}
$L_2(12,4)\geq6$.
\end{proposition}

\begin{proof}
The six distinct words in \eqref{eq:codebook} form a four-deletion-correcting
code by the table above.
\end{proof}

The released checker confirms this claim two ways.  It enumerates the
four-deletion descendants of each word by keep/delete recursion, and it
computes all 15 LCS values by dynamic programming.  The six descendant sets
contain 198 distinct words in total, with no collision between codewords.

\section{A weighted certificate lemma}

The upper bound puts nonnegative weights on two families of packing
inequalities.  We state the argument in the form the certificate uses.

\begin{lemma}[Weighted descendant-and-clique bound]\label{lem:weighted}
Fix $q,n,t$.  Let $a_d\geq0$ for $d\in[q]^{n-t}$.  Let
$K_1,\dots,K_m$ be cliques in the conflict graph, with weights $b_i\geq0$.
Suppose each $x\in[q]^n$ satisfies
\begin{equation}\label{eq:coverage}
  \sum_{d\in\D_t(x)}a_d+\sum_{i:x\in K_i}b_i\geq1.
\end{equation}
Then each $t$-deletion-correcting code $\C\subseteq[q]^n$ satisfies
\begin{equation}\label{eq:weightedbound}
  |\C|\leq \sum_{d\in[q]^{n-t}}a_d+\sum_{i=1}^m b_i.
\end{equation}
\end{lemma}

\begin{proof}
Sum \eqref{eq:coverage} over $x\in\C$ and exchange the order of summation:
\begin{align*}
|\C|
&\leq \sum_{x\in\C}\left(
       \sum_{d\in\D_t(x)}a_d+\sum_{i:x\in K_i}b_i\right)\\
&=\sum_d a_d|\C\cap R_d|+\sum_i b_i|\C\cap K_i|.
\end{align*}
The code meets each row $R_d$ and each clique $K_i$ in at most one word.
Nonnegativity of the weights gives \eqref{eq:weightedbound}.
\end{proof}

The clique terms carry the proof.  With weights on the 256 descendant rows
alone, the best bound this lemma can give is the optimum of the fractional
packing LP, which HiGHS reports as about $7.054$.  That yields only
$L_2(12,4)\leq7$, the same bound the variable-length table implies.  Adding 20 cliques pushes the bound below seven,
and each clique remains easy to check: every pair of its words must share a
length-eight subsequence.

\section{The exact certificate}\label{sec:certificate}

Set $M=1{,}000{,}000$.  The certificate stores integers $A_d\geq0$ for
$d\in\bits^8$ and $B_i\geq0$ for 20 conflict cliques $K_i$, and sets
\[
  a_d=A_d/M,\qquad b_i=B_i/M.
\]
The JSON file omits zero descendant weights.  The checker verifies the
following facts in exact arithmetic.

\begin{center}
\begin{tabular}{@{}lr@{}}
\toprule
Quantity & Verified value\\
\midrule
Length-12 words checked & 4,096\\
Possible length-8 descendants & 256\\
Nonzero descendant weights & 117\\
Conflict cliques & 20\\
Clique sizes & 62 to 94\\
Clique pairs checked & 61,263\\
Minimum integer word coverage & 1,000,000\\
Sum of all integer weights & 6,999,616\\
\bottomrule
\end{tabular}
\end{center}

So for each $x\in\bits^{12}$,
\begin{equation}\label{eq:integercoverage}
  \sum_{d\in\D_4(x)}A_d+\sum_{i:x\in K_i}B_i\geq M,
\end{equation}
while
\begin{equation}\label{eq:integerobjective}
  \sum_d A_d+\sum_i B_i=6{,}999{,}616<7M.
\end{equation}

\begin{proposition}\label{prop:upper}
$L_2(12,4)\leq6$.
\end{proposition}

\begin{proof}
Divide \eqref{eq:integercoverage} and \eqref{eq:integerobjective} by $M$ and
apply \cref{lem:weighted}.  Each valid code $\C$ satisfies
\[
  |\C|\leq\frac{6{,}999{,}616}{1{,}000{,}000}=6.999616<7.
\]
Since $|\C|$ is an integer, $|\C|\leq6$.
\end{proof}

\Cref{prop:lower,prop:upper} together prove the main result.

\begin{theorem}\label{thm:main}
In the ordinary ordered binary deletion model,
\[
  \boxed{L_2(12,4)=6}.
\]
\end{theorem}

\section{Finding and checking the certificate}

We separate the search, which no one needs to trust, from the small
computation that the proof relies on.

\subsection{Search}

We treated the 4,096 words of length twelve as vertices of the conflict graph
and solved the linear-programming relaxation of the 256 descendant-row packing
constraints with HiGHS.  While the LP optimum stayed at seven or above, we
solved an integer maximum-weight clique problem on the support of the current
LP solution to find a violated clique inequality, then added it as a cut.
After 20 cuts the optimum fell below seven.  We multiplied the nonnegative dual
multipliers by $M$, rounded each one up, and kept the rounded vector only after
the exact coverage and total checks passed.  \Cref{thm:main} uses none of the
floating-point output.

\subsection{Standalone checker}

The verifier is a Python program that imports only the standard library.
Given the JSON certificate, it checks that:
\begin{enumerate}[leftmargin=2.2em]
  \item the declared model and parameters equal $(q,n,t)=(2,12,4)$;
  \item each word and descendant is a fixed-width binary string;
  \item each clique lists distinct words, and all 61,263 within-clique pairs
        share a four-deletion descendant, reconstructed from scratch;
  \item each of the 4,096 words satisfies \eqref{eq:integercoverage};
  \item the stored total equals 6,999,616, below $7M$; and
  \item the certificate bytes match the published SHA-256 digest.
\end{enumerate}

A separate program, independent of the search code, checks the six-word code
through both descendant-set disjointness and the LCS criterion.  Mutation
tests lower a necessary weight and corrupt a clique, and the verifier rejects
both.

\begin{center}
\begin{tabular}{@{}p{0.28\linewidth}p{0.64\linewidth}@{}}
\toprule
Artifact & SHA-256\\
\midrule
Codebook &
\texttt{\footnotesize 5dee9d5e15202ebe6ecff0fe0f35e704}\newline
\texttt{\footnotesize e3b8d957967a9116becdeaa6cc1c023f}\\[0.4em]
Upper certificate &
\texttt{\footnotesize d9c8b720d112957250c3c06ec786634a}\newline
\texttt{\footnotesize 573d6e1f78a449f2b2da5a829cf7fc82}\\
\bottomrule
\end{tabular}
\end{center}

\section{Lean 4 formalization}

The Lean development defines the ordered deletion model, imports the
certificate as transparent data, checks its finite facts, and proves the
weighted-sum argument.  Words of length twelve are elements of
\lstinline|Fin 4096| and descendants are elements of \lstinline|Fin 256|.
The predicate \lstinline|hasDescendant| compares fixed-width bit lists with
\lstinline|List.isSublist|, which tests for a subsequence.  A family
\lstinline|code : Fin k -> Word| is \lstinline|ValidFamily| when it is injective
and no two of its words share a length-eight descendant.

The main theorem reads:
\begin{lstlisting}[language=Lean]
theorem binary_length12_four_deletion_exact :
    (lowerCode.Nodup /\ lowerCode.length = 6 /\
      explicitCodeValid lowerCode = true) /\
    Not (Exists fun code : Fin 7 -> Word => ValidFamily code)
\end{lstlisting}

The second conjunct excludes every injective labelled family of seven
compatible words.  A code with seven or more words contains such a family, so
the two conjuncts together state \cref{thm:main}.  The proof isolates four
finite checks from the structural argument:
\begin{itemize}[leftmargin=2.2em]
  \item validity and size of the six-word code;
  \item validity of the 20 conflict cliques;
  \item coverage of all 4,096 words; and
  \item the exact total-weight inequality.
\end{itemize}

The development contains no \lstinline|sorry|, no \lstinline|admit|, and no
project-local axiom.  Its four finite checks use \lstinline|native_decide|, so
the Lean proof trusts Lean's compiler and runtime in addition to its kernel;
\lstinline|#print axioms| lists exactly these four leaves beyond the standard
axioms.  Replacing them with kernel-checked proof terms would remove that
dependence.  The Python checker does not share it.

\section{Reproducibility}

From the \texttt{formalization} directory of the accompanying repository, run:
\begin{lstlisting}
python3 research/deletion_code_attempt/verify_code.py
python3 research/deletion_code_attempt/proof/verify_upper_certificate.py
python3 research/deletion_code_attempt/proof/test_upper_certificate_mutations.py
lake build DeletionCodes.Conjectures.DeletionCode124
lake build
\end{lstlisting}

The Python programs end with these markers:
\begin{itemize}[leftmargin=2.2em]
  \item \texttt{VERIFICATION=PASS} for the six-word code;
  \item \texttt{UPPER\_CERTIFICATE=PASS} for the upper certificate; and
  \item \texttt{MUTATION\_AUDIT=PASS} for the mutation tests.
\end{itemize}
The Lean build prints the axioms of the main theorem, so a reader can inspect
its trust boundary.

\section{Concluding remarks}

This paper settles one finite entry, not an asymptotic question.  A reader can
check the whole proof: six lines give the construction, one weighted sum gives
the upper bound, and a versioned certificate holds every finite fact, checked
in both Python and Lean.  The same combination of descendant rows and clique
cuts applies without change to other open entries.  The nearest are
$6\leq L_2(13,4)\leq14$ and $7\leq L_2(14,4)\leq22$ in Kolev's table, which
lie beyond the range of the variable-length tables.

\paragraph{Literature search.}
Searches through September 28, 2026 found the bounds in
\cref{tab:history} and no source that states the exact value or an upper bound
below seven.  Recent finite-length studies stop at three deletions, and
arXiv listings from July through September 2026 contain no competing
result.  We could not obtain the 2015 thesis that Kolev cites,
and an unpublished result could exist.  The theorem and its certificate stand
independent of any priority question.

\paragraph{Use of AI tools.}
The author used AI coding and reasoning assistants to write search code, the
certificate checkers, the Lean formalization, and drafts of this manuscript.
The author takes responsibility for the statements, the literature claims, the
released artifacts, and the correctness of the paper.

\bibliographystyle{alpha}
\bibliography{references}

\end{document}
